% Regression: t2_idempotent — imported current-behavior fixture from the handoff corpus.
% Target: RMP 2011.12127 fig3_fig3-pepe_Idempotent.pdf (the closed MPO
% ring around one PEPS site is an idempotent).
% Formula:  ring_j( T_i ) = delta_{i,j} T_i  —  the virtual MPO string
% (red) closed into a ring through four MPO tensors around one PEPS
% tensor T_i (blue, four virtual legs + one physical leg), containing one
% marked MPO tensor j (blue parallelogram) whose free string end exits;
% closing the ring projects: it equals delta_{i,j} times the bare tensor
% with one residual string leg.
% Genuinely non-grid (a ring threading glyphs around a star): tenkzfree.
% Honest gaps: \tnjoin has no wire-hue key (the red MPO ring vs black
% lattice wires IS the two-level structure) and \tnput has no marked/hue
% skin (the blue T_i and blue j render as house ink); the mpo skin is a
% house box, not the paper's slanted parallelogram.
\documentclass[varwidth=19cm, border=6pt]{standalone}
\usepackage{amsmath, amssymb}
\usepackage{tenkz}
\begin{document}

\[
  \begin{tenkzfree}
    % the PEPS tensor: physical leg up, horizontal + diagonal lattice
    % wires; the label takes the free south side (west carries the wire)
    \tnput[dot, label pos=south]{C}{(0,0)}{i}
    \tnjoin{C.north}{0mm,7mm}
    % four MPO tensors of the ring, on the lattice wires
    \tnput[mpo]{W}{(-16mm,0)}{}
    \tnput[mpo]{E}{(16mm,0)}{}
    \tnput[mpo]{N}{(8mm,10mm)}{}
    \tnput[mpo]{S}{(-8mm,-10mm)}{}
    % horizontal wire through W, C, E
    \tnjoin{-24mm,0mm}{W.west}
    \tnjoin{W.east}{C.west}
    \tnjoin{C.east}{E.west}
    \tnjoin{E.east}{24mm,0mm}
    % diagonal wire through S, C, N
    \tnjoin{-14mm,-17mm}{S.south west}
    \tnjoin{S.north east}{C.south west}
    \tnjoin{C.north east}{N.south west}
    \tnjoin{N.north east}{14mm,17mm}
    % the marked MPO tensor j on the ring, free string end exiting SE
    \tnput[mpo]{J}{(12mm,-9mm)}{j}
    \tnjoin{J.south east}{20mm,-15mm}
    % the MPO ring: arcs threading W -> N -> E -> J -> S -> W
    \tnjoin[route=curve, out=90, in=180]{W.north}{N.west}
    \tnjoin[route=curve, out=0, in=90]{N.east}{E.north}
    \tnjoin[route=curve, out=-90, in=45]{E.south}{J.north east}
    \tnjoin[route=curve, out=225, in=0]{J.south west}{S.east}
    \tnjoin[route=curve, out=180, in=-90]{S.west}{W.south}
  \end{tenkzfree}
  \;=\;\delta_{i,j}\;
  \begin{tenkzfree}
    % the bare tensor: same legs, one residual string stub to the SE
    \tnput[dot, label pos=south]{D}{(0,0)}{i}
    \tnjoin{D.north}{0mm,7mm}
    \tnjoin{-10mm,0mm}{D.west}
    \tnjoin{D.east}{10mm,0mm}
    \tnjoin{-7mm,-9mm}{D.south west}
    \tnjoin{D.north east}{7mm,9mm}
    \tnjoin{D.south east}{8mm,-6mm}
  \end{tenkzfree}
\]

\end{document}
